% Generated from locale/tr/content/proof-theory/propositions-as-types/normalization.tex; do not hand-edit.


\olfileid[tr]{int}{pty}{nor}

\olsection{Normalleştirme}

Bu kısımda tipli ispat terimleri için \emph{$\beta$-normalleştirme}
sonucunu ele alıyoruz. Tip olarak önerme karşılığından yararlanarak her ispat
teriminin, aşağıda ölçülen üç $\beta$-indirgenen ailesi bakımından normal bir
biçime indirgenebildiğini göstereceğiz; doğal tümdengelim !!{derivation}s için
karşılık gelen $\beta$-indirgemeleri hakkındaki sonuç da bundan çıkacaktır.
Yer değiştirme dönüşümlerini de kapsayan tam normalleştirme savı, bu ölçünün
kapsamı dışındadır.

Önce terimlerin karmaşıklığını ölçen bazı fonksiyonlar tanımlayalım.
!!a{formula}s için \emph{uzunluk}~$\len{!A}$ şöyle tanımlanır:
\begin{align*}
  \len{p} &= 0\\
  \len{!A \land !B} &= \len{!A} + \len{!B} + 1\\
  \len{!A \lor !B} &= \len{!A} + \len{!B} + 1 \\
  \len{!A \lif !B} &= \len{!A} + \len{!B} + 1.
\end{align*}
Bir indirgenen~$M$'nin karmaşıklığı, onun \emph{kesme
  derecesi}~$\cutrank{M}$ ile ölçülür:
\begin{align*}
  \cutrank{(\lambd[\typeof{x}{!A}][\typeof{N}{!B}])Q} &=
  \len{!A} + \len{!B} + 1 \\
  \cutrank{\ande{i}{\andi{\typeof{M}{!A}}{\typeof{N}{!B}}}} &=
  \len{!A} + \len{!B} + 1 \\
  \cutrank{\ore{\ori{i}{!A_{i}}{\typeof{M}{!A_i}}}
    {\typeof{x_1}{!A_1}}{\typeof{N_1}{!C}}
    {\typeof{x_2}{!A_2}}{\typeof{N_2}{!C}}} &=
  \len{!A_1} + \len{!A_2} + 1.
\end{align*}
Bir ispat teriminin $\beta$-karmaşıklığı, içerdiği en karmaşık
$\beta$-indirgenenin kesme derecesidir; terim $\beta$-normalse bu değer~$0$'dır:
\[
  \maxrank{M} = \max \bigl(\{\cutrank{N} \mid N, M\text{'nin indirgenen
    bir alt terimidir}\}\cup\{0\}\bigr).
\]

\begin{lem}\ollabel{lem:subst}
  $\Subst{M}{\typeof{N}{!A}}{\typeof{x}{!A}}$ bir indirgenense ve $M
  \not\ident x$ ise aşağıdaki durumlardan biri geçerlidir:
  \begin{enumerate}
  \item $M$'nin kendisi bir indirgenendir;
  \item $M$, $\ande{i}{x}$ biçimindedir ve $N$, $\andi{P_1}{P_2}$
    biçimindedir;
  \item $M$, $\ore{i}{x_1}{P_1}{x_2}{P_2}$ biçimindedir ve $N$,
    $\ori{i}{}{Q}$ biçimindedir;
  \item $M$, $x Q$ biçimindedir ve $N$, $\lambd[x][P]$ biçimindedir.
  \end{enumerate}
  İlk durumda $\cutrank{\Subst{M}{N}{x}} = \cutrank{M}$; öteki
  durumlarda $\cutrank{\Subst{M}{N}{x}} = \len{!A}$ olur.
\end{lem}
\begin{proof}
  $M$ üzerinde tümevarımla.
  \begin{enumerate}
  \item $M$ tek bir $y$ değişkeniyse ve $y \not\ident x$ ise
    $\Subst{y}{N}{x}$, $y$'dir; dolayısıyla indirgenen değildir.
  \item $M$, $\andi{N_1}{N_2}$, $\lambd[x][N]$ ya da
    $\ori{i}{!A}{N}$ biçimindeyse
    $\Subst{M}{\typeof{N}{!A}}{\typeof{x}{!A}}$ da aynı biçimdedir ve
    dolayısıyla indirgenen değildir.
  \item $M$, $\ande{i}{P}$ biçimindeyse iki durumu ele alırız.
    \begin{enumerate}
      \item $P$, $\andi{P_1}{P_2}$ biçimindeyse $M \ident
        \ande{i}{\andi{P_1}{P_2}}$ bir indirgenendir ve açıkça
        \[
        \Subst{M}{N}{x} \ident
        \ande{i}{\andi{\Subst{P_1}{N}{x}}{\Subst{P_2}{N}{x}}}
        \]
        de bir indirgenendir. Kesme dereceleri eşittir.
      \item $P$ tek bir değişkense yerine koymanın bir indirgenen üretmesi
        için bu değişken $x$ olmalı, $N$ de $\andi{P_1}{P_2}$ biçiminde
        olmalıdır. Şimdi
        \[
        \Subst{M}{N}{x} \ident \Subst{\ande{i}{x}}{\andi{P_1}{P_2}}{x}
        \]
        terimini ele alalım. Bu terim
        $\ande{i}{\andi{P_1}{P_2}}$'dir ve kesme derecesi
        $\len{!A}$'dır.
    \end{enumerate}
  \end{enumerate}
  $\ore{N}{x_1}{N_1}{x_2}{N_2}$ ve $P Q$ durumları benzerdir.
\end{proof}

\begin{lem}
  $M$, $M'$'ye büzülüyorsa ve $M$'nin indirgenen olan her öz alt
  terimi~$N$ için $\cutrank{M} > \cutrank{N}$ ise
  $\cutrank{M} > \maxrank{M'}$ olur.
\end{lem}

\begin{proof}
  Durumlara göre ispatlayalım.
  \begin{enumerate}
  \item $M$, $\ande{i}{\andi{M_1}{M_2}}$ biçimindeyse $M'$, $M_i$'dir.
    $M_i$'nin her alt terimi aynı zamanda $M$'nin öz alt terimi olduğundan
    iddia geçerlidir.
  \item $M$,
    $(\lambd[\typeof{x}{!A}][N])\typeof{Q}{!A}$ biçimindeyse $M'$,
    $\Subst{N}{\typeof{Q}{!A}}{\typeof{x}{!A}}$'dır. $M'$ içindeki bir
    indirgeneni ele alalım. Ya $N$ içinde aynı kesme derecesine sahip
    karşılık gelen bir indirgenen vardır ve varsayım gereği bu derece
    $\cutrank{M}$'den küçüktür; ya da kesme derecesi $\len{!A}$'dır ve bu
    da tanım gereği
    $\cutrank{(\lambd[x^{!A}][N])Q}$'dan küçüktür.
  \item $M$,
    \[
    \ore{\ori{i}{}{\typeof{N}{!A_i}}}
        {\typeof{x_1}{!A_1}}{\typeof{N_1}{!C}}
        {\typeof{x_2}{!A_2}}{\typeof{N_2}{!C}}
    \]
    biçimindeyse $M' \ident \Subst{N_i}{N}{\typeof{x_i}{!A_i}}$ olur.
    $M'$ içindeki bir indirgeneni ele alalım. Ya $N_i$ içinde aynı kesme
    derecesine sahip karşılık gelen bir indirgenen vardır ve varsayım gereği
    bu derece $\cutrank{M}$'den küçüktür; ya da kesme derecesi
    $\len{!A_i}$'dır ve bu da tanım gereği
    $\cutrank{\ore{\ori{i}{}{\typeof{N}{!A_i}}}
      {\typeof{x_1}{!A_1}}{\typeof{N_1}{!C}}
      {\typeof{x_2}{!A_2}}{\typeof{N_2}{!C}}}$'dan küçüktür.
  \end{enumerate}
\end{proof}

\begin{thm}
  Bütün ispat terimleri, yukarıdaki üç $\beta$-indirgenen ailesi bakımından
  normal bir biçime indirgenir; bütün !!{derivation}s da karşılık gelen
  $\beta$-indirgemeler bakımından normal !!{derivation}s elde edilecek biçimde
  indirgenir.
\end{thm}

\begin{proof}
  İkinci iddia birinciden çıkar. Birinciyi, bir ispat terimi~$M$ için
  $m=\maxrank{M}$ üzerinde güçlü tümevarımla ispatlayalım.
  \begin{enumerate}
  \item $m=0$ ise $M$ zaten bu anlamda $\beta$-normaldir.
  \item Aksi durumda, $M$ içinde kesme derecesi~$m$ olan indirgenenlerin
    sayısı~$n$ üzerinde tümevarım yaparız.
    \begin{enumerate}
    \item $n=1$ ise $m = \cutrank{N}$ olacak ve indirgenen olan her öz
      alt terim~$P$ için $\cutrank{N} > \cutrank{P}$ olacak biçimde bir
      indirgenen~$N$ seçin. Böyle bir indirgenen vardır; çünkü her terimin
      yalnızca sonlu sayıda alt terimi bulunur.

      $N$'nin indirgenmişini $N'$ ile gösterelim. Önceki lemma gereği
      $\maxrank{N'} < \maxrank{N}$ olur. Böylece kesme derecesi~$m$ olan
      indirgenenlerin sayısı sıfıra iner; dolayısıyla bütün terimin azami
      kesme derecesi azalır ve $m$ için tümevarım varsayımını uygulayabiliriz.
    \item Tümevarım adımı için $n>1$ varsayın. İşlem benzerdir: kesme
      derecesi~$m$ olan ve aynı derecede indirgenen bir öz alt terim
      içermeyen bir~$N$ seçip büzün. Bu kez $n$ yalnızca daha küçük pozitif
      bir sayıya iner ve $m$ değişmez. $n$ için tümevarım varsayımını
      uygularız.
    \end{enumerate}
  \end{enumerate}
\end{proof}

Tipli $\lambd$-terimlerinin $\beta$-normalleşmesi, böyle her terimin bir
$\beta$-normal biçimi olduğu anlamına gelir. $\lambd$-terimlerini $\beta$-indirgemeyle
normal biçimlerine getirmeyi bir hesaplamayı yürütmek, normal biçimi de
hesaplamanın değeri olarak düşünürsek her tipli $\lambd$-terimi için sonunda
duran ve bir değer üreten bir hesaplama yolu vardır. Ayrıca tip korunumu, değerin
doğru tipte olmasını sağlar. Örneğin $N$'nin tipi $!A \lif !B$ ise (yani
$N$, tipi~$!A$ olan argümanlar üzerinde çalışıp tipi~$!B$ olan değerler
üretmesi gereken bir fonksiyonsa) ve onu tipi~$!A$ olan bir $M$ terimine
(girdiye) uygularsak $(NM)$, bir $N'$ terimine (çıktıya) indirgenir ve bu
çıktının tipi~$!B$ olmak zorundadır.

Standart tipli $\lambd$-kalkülüsü için daha güçlü bir sonuç da geçerlidir:
alt terimlere $\beta$-indirgemeleri uygulamanın yalnızca belirli bir yolunda
değil, \emph{her} yolunda sonunda normal biçime ulaşılır. Bu özelliğe
\emph{güçlü normalleştirme} denir. Kaynak metin bu güçlü sonucu burada
ispatlamaz; aşağıdaki Church--Rosser tartışmasında onu açık bir varsayım olarak
kullanır.

Bir hesaplamayı yürütmenin farklı yollarının \emph{aynı} değeri ürettiğini
nasıl biliriz? Şimdiye kadar tip korunumundan yalnızca bütün sonuçların aynı
tipte olduğunu biliyoruz. Tipli $\lambd$-kalkülüsünün başka bir önemli
özelliği daha vardır: sistem \emph{Church--Rosser özelliğine sahiptir}.
$N \red N_1$ ve $N \red N_2$ ise $N_1 \red N'$ ve
$N_2 \red N'$ olacak biçimde bir $N'$ terimi vardır. $N \red N'$ ifadesinin,
$N'$'nin sıfır ya da daha çok $\redone$-indirgeme adımıyla $N$'den elde
edildiği anlamına geldiğini anımsayın. $N_1$ normal biçimdeyse $N_1 \red N'$
olacak tek $N'$ terimi, sıfır $\redone$-indirgeme adımıyla elde edilen
$N_1$'in kendisidir. Dolayısıyla hem $N_1$ hem $N_2$ normal biçimdeyse
$N_1=N'=N_2$ olur. Başka bir deyişle, $\red$ güçlü normalleştirici olduğu
için bir terimi indirgemenin her yolu sonunda normal biçime ulaşır;
$\red$ Church--Rosser özelliğine sahip olduğu için de her iki normal biçim
eşittir. Böylece tipli $\lambd$-kalkülüsünde hesaplama her zaman sona erer ve
hesaplamanın nasıl ilerlediğinden bağımsız olarak aynı değeri (normal biçimi)
verir.

Bir bağıntının Church--Rosser özelliğine sahip olduğunu ispatlamak genellikle
güçtür. Ancak şu daha
zayıf koşulu kolayca gösterebiliriz:

\begin{prop}[Zayıf Church--Rosser özelliği (ispat taslağı)]
$N \redone N_1$ ve $N \redone N_2$ ise $N_1 \red N'$ ve $N_2 \red N'$
olacak biçimde bir $N'$ terimi vardır.
\end{prop}

\begin{proof}
  Kaynak metnin taslağı şu durumları ayırır. $N_1$ ve $N_2$, sırasıyla $N$
  içindeki indirgenen $M_1$ ve $M_2$'nin
  indirgenmişleriyle değiştirilmesiyle elde edilir. $M_1$ ile $M_2$, $N$'nin
  alt terimleridir. Aynı indirgenen seçilmişse $N_1=N_2$ olur ve sonuç
  dolaysızdır. Aksi hâlde bunlar ya çakışmayan farklı konumlarda bulunur ya
  da biri ötekinin alt terimidir. İlk durumu ele alalım: $N$, $N(M_1,M_2)$
  biçimindedir. $M_1$'i indirgenmişi $M_1'$ ile değiştirmek
  $N(M_1',M_2)$ olan $N_1$'i; $M_2$'yi indirgenmişi $M_2'$ ile değiştirmek
  ise $N(M_1,M_2')$ olan $N_2$'yi verir. Her ikisi de
  $N(M_1',M_2')$'ye indirgenir. Öteki durumda $M_2$'nin $M_1$'in alt terimi
  olduğunu varsayın (ters durum simetriktir) ve $M_1$'in indirgenen biçimine
  göre durumları ayırın. Örneğin
  $M_1 \ident \proj{1}{\pair{O_1}{O_2}}$ ise $O_1$'e indirgenir. $M_2$,
  $O_1$'in bir alt terimiyse $M_2$'yi indirgemek $O_1'$i verir. Böylece önce
  $M_1$'i, ardından $M_2$'yi indirgemek $O_1'$e ulaşır. Öte yandan önce
  $M_2$'yi indirgemek $\proj{1}{\pair{O_1'}{O_2}}$ terimini verir; bu terim
  de $O_1'$e indirgenir. Öteki indirgenen biçimleri için de ayrı bir kritik
  çift denetimi gerekir; kaynak metin bu sonlu denetimi ayrıntılandırmaz.
\end{proof}

\begin{lem}[Newman lemması]
$\red$ güçlü normalleştiriciyse ve zayıf Church--Rosser özelliğine sahipse
Church--Rosser özelliğine sahiptir.
\end{lem}

\begin{proof}
  $\red$ güçlü normalleştirici olduğundan her indirgeme dizisi bir normal
  biçimde sona erer. Bu koşul altında normal biçimlerin tekliği,
  $\red$'in Church--Rosser özelliğine sahip olmasına denktir. Bir $N$
  teriminin, şu indirgeme
  dizileriyle elde edilen iki farklı $N'$ ve $N''$ normal biçimi bulunduğunu
  varsayalım:
  \begin{align*}
  N &\redone N_1' \redone N_2' \redone \dots \redone N_k' = N'
    \text{ ve}\\
  N &\redone N_1'' \redone N_2'' \redone \dots \redone N_l'' = N''.
  \end{align*}
  (İndirgeme dizileri sıfır uzunluklu olamaz; yoksa $N$ zaten normal biçimde
  olurdu.)

  $N_1'=N_1''$ ise bu ortak terimin iki farklı normal biçimi ($N'$ ve
  $N''$) vardır. $N_1'\neq N_1''$ ise zayıf Church--Rosser özelliği gereği
  $N_1'\red N'''$ ve $N_1''\red N'''$ olacak biçimde bir $N'''$ vardır.
  Güçlü normalleştirme gereği $N'''$ bir normal biçim~$P$'ye indirgenir.
  $N'\neq N''$ olduğundan $P\neq N'$ ya da $P\neq N''$ olur. Dolayısıyla
  $N_1'$ ile $N_1''$ terimlerinden en az birinin iki farklı normal biçimi
  vardır. Böylece $N$'nin iki farklı normal biçimi varsa
  $N\redone N_1$ olacak ve $N_1$'in de iki farklı normal biçimi bulunacak
  bir $N_1$ teriminin var olduğunu gösterdik. Aynı akıl yürütmeyi sürdürerek
  $N\redone N_1\redone N_2\redone\dots$ sonsuz dizisini elde ederiz. Ancak
  bu, $\redone$ güçlü normalleştirici olduğu için olanaksızdır.
\end{proof}

